\documentclass[12pt]{amsart}
\begin{document}\bigskip\bigskip\bigskip\bigskip
\centerline{\bf MATHEMATICS 6130}
\centerline{Homework 1}
\begin{enumerate}
\item Define curve, irreducible curve and rational curve.

\item Show that if $k$ is algebraically closed, then any curve has infinitely many points.

\item Show that $y^2=x^2(x+1)$, $x^2+2xy=3$ and $y^2=x^3$ are rational (over an algebraically closed field).

\item Show that (if $k$ is algebraically closed) an irreducible curve determines its equations (up to units). (You may assume the following lemma: If $f\in k[x,y] $ is irreducible and does not divide $g\in k[x,y] $, then the set of common solutions is finite.)

\item State Luroth's Theorem.

\item Define the field of rational functions $k(X)$ of a plane curve $X$. When is a rational function regular at a point $p\in X$?

\item Show that a plane curve $X$ is rational if and only if $k(X)\cong k(t)$ ($k=\bar k$).

\item Define rational map between two plane curves.
\end{enumerate}
\centerline{Homework 2}
\begin{enumerate}
\item Define nonsingular points, point of multiplicity $r$, cusps, node of a plane curve.
\item What can you say about a plane curve of degree $d$ with a point of multiplicity $d,d-1,d-2$?
\item Show that irreducible plane curves have finitely many singular points.
\item Define local parameter of an irreducible plane curve.
\item Define the tangent line and the flex of a curve at a nonsingular point.
\item Define $\mathbb P ^n$. Explain why polynomial functions are not defined on $\mathbb P ^n$, but the zeroes of homogeneous polynomials are defined.
\item Show that all circles contain the point at infinity $(1,i,0)$ and $(1,-i,0)$.
\item Show that a parabola is tangent to the line at infinity (eg. $y=x^2$ is tangent at the point $[0:1:0]$).
\item Which rational functions in $k(x)$ are regular at the point at infinity $[1:0]\in \mathbb P ^1$?
\item Show that if $f\in k(x)$ is regular at every point in $\mathbb P ^1$, then $f$ is constant.
\item Prove that an irreducible cubic has at most 1 singular point.
\item Show that if the characteristic of $k$ is $p>0$, then every line through the origin is tangent to $y=x^{p+1}$.

\end{enumerate}
\centerline{Homework 3}
\begin{enumerate}
\item Define closed subsets of $\mathbb A ^n$ and show that they define a topology.
\item Define the ring of rational functions $k[X]$ of a closed subset $X\subset \mathbb A \
^n$ and a regular map between closed subsets. When are two closed subsets isomorphic?
\item Show how to recover a closed subset from
its ring of regular function and the rational map $X\to Y$
between two closed subsets from a homomorphism between their rings of rational functions $\
k[Y]\to k[X]$.
\item State the Nullstellensatz.
\item Define the Frobenius map $\phi :\mathbb   A ^n\to \mathbb A ^n$ and show that if $X\\
subset\
 \mathbb A ^n$ is defined over $\mathbb F_p$ (i.e. it's ideal is generated by functions in\
 $\mathbb     F_p[x_1,\ldots,x_n ]$) then $\phi$ defines a rational map from $X$ to $X$.
\item Let $I_X$ be the ideal of a closed subset $X\subset\
 \mathbb A ^n$. Show that $k[X]=k[\mathbb A ^n]/I_X$ has no nilpotents.
\end{enumerate}
\centerline{Homework 4}
\begin{enumerate}
\item Let $X\subset \mathbb A ^n$ be an irreducible closed subset.Define the field of rational functions $k(X)$. When is $\phi\in k(X)$ regular at $p\in X$? Show that the domain of definition of  $\phi\in k(X)$ is a non-empty open subset.
\item Let $X\subset \mathbb A ^n$ be an irreducible closed subset. Show that if $\phi \in k(X)$ is regular at all $x\in X$, then $\phi \in k[X]$.
\item Define rational maps $\phi:X\to Y$ where $X\subset \mathbb A ^n$ and $Y\subset \mathbb A^m$ are irreducible closed subsets. When is $\phi$ dominant?
Prove that $\phi $ is dominant iff $\phi ^*:k[Y]\to k(X)$ is injective.
\item Define birational map and show that a map is birational iff it induces an isomorphism of function fields.
\item Show that if $I\subset k[x_0,\ldots , x_n]$ is a homogeneous ideal defining the empty subset of $\mathbb P ^n$, then $I\supset m^s$ where $m=(x_0,\ldots , x_n)$.\end{enumerate}
\centerline{Homework 5}
\begin{enumerate}
\item Define quasi-projective variety, regular functions on quasi-projective varieties and rational map / regular map between quasi-projective varieties.
\item Define principal affine open subsets and show that any quasi-projective variety can be covered by  principal affine open subsets.
\item Explain why a regular map of quasi-projective varieties is continuous.
\item Define $\mathcal O _X$ and $k(X)$ (the function field of $X$) for an irreducible quasi-projective variety $X$.
\item Show that if $X$ is an irreducible quasi-projective variety then $k(X)=k(\bar X)$ (however $k[X]\ne k[\bar X]$ usually) and if $X$ is an irreducible affine variety then . $k(X)$ is the field of rational functions of $X$.
\item Show that $\mathbb A ^2\setminus (0,0)$ is not affine.
\item Define the projection from a linear subspace $E\subset \mathbb P ^n$ and the Veronese embedding $\nu _m:\mathbb P^n\to \mathbb P ^N$ where $N=\binom {n+m}n -1$.
\end{enumerate}
\centerline{Homework 6}
\begin{enumerate}
\item Define a closed embedding $\phi : \mathbb P ^n\times \mathbb P ^m\to \mathbb P^N$ and show that $\phi ( \mathbb P ^n\times \mathbb P ^m)\subset  \mathbb P^N$ is closed. Explain why $\phi$ is closed.
\item Explain why a regular map from a projective variety to an affine variety is constant (you may use the fact that the image of a projective variety under a regular map is closed).
\item Explain why the set of reducible homogeneous polynomials of degree $m$ in $n+1$ variables is a proper closed subset of the set of all homogeneous polynomials of degree $m$ parametrized by $\mathbb P ^N$ where $N =\binom{m+m}m -1$.
\item Define finite map of affine varieties. Show that all fibers of any such map have finite cardinality but the converse does not hold.
\item Show that a finite map is surjective (you may use the fact that if $B$ is a finite $A$ module, $1_B\in A\subset B$ a subring, then for any proper ideal $(1)\ne a\subset A$ we have that $aB\ne B$).
\item Prove that $\mathbb A ^2\setminus x$, $\mathbb P ^2\setminus x$, and 
$\mathbb P^1\times \mathbb A ^1$ are neither affine nor projective varieties.  
\end{enumerate}

\centerline{Homework 7}
\begin{enumerate}
\item Prove that if $X\subset \mathbb P ^N$ is a projective variety and $E\subset \mathbb P ^N$ is a $d$-dimensional linear subspace such that $E\cap X=\emptyset$, then the projection $\pi : \mathbb P ^N\to \mathbb P ^{N-d-1}$ is finite on to its image.
\item Deduce from the previous exercise that if $F_0, \ldots , F_s$ are homogeneous of degree $m$ and $X\subset \mathbb P ^N$ is a projective variety such that $X\cap V(F_0, \ldots , F_s)=\emptyset $, then the rational map $\phi :X\dasharrow \mathbb P ^s$ defined by $ (F_0, \ldots , F_s)$ is finite.
\item Let $X$ be a quasi-projective variety. Define $\dim X$ and show that if $X$ is irreducible then $\dim X$ is a birational invariant.
\item Show that if $X\subset Y$ is an inclusion of quasi projective varieties, then $\dim X\leq \dim Y$.
\item Show that every irreducible component of a hyprsurface in $\mathbb A ^N$ has codimension $1$.
\item Give alternative definitions of $\dim X$ in terms of the Noether Normalization theorem and by sequences of irreducible proper closed subsets.
\item Let $X,Y\subset \mathbb P ^N$ be irreducible quasi-projective varieties. Show that $codim (X\cap Y)\leq codim (X)+codim (Y)$.
\item Let $f:X\to Y$ be a regular map of  irreducible quasi-projective varieties. Show that each (non-empty) fiber has dimension $\geq \dim X-\dim Y$. What can you say about the function $\dim (f^{-1}(y)$ where $y\in Y$? 
\end{enumerate}
\centerline{Homework 8+9}
\begin{enumerate}
\item Let  $x\in X$ be a point on an irreducible affine variety. Define $\mathcal O _{x,X}$.
\item Let  $x\in X$ be a point on an irreducible affine variety. Define $\Theta _{x,X}$ and show that $\Theta _{x,X}\cong (\frak m _x/\frak m_x ^2)^\vee$.
\item Prove that the curve $C\subset \mathbb A ^N$ given by $t\to (t^N,t^{N+1},\ldots , t ^{2N-1})$ can not be embedded in $\mathbb A ^n$ for any $n<N$.
\item Explain why there exists an open subset $X_{sm}\subset X$ such that $\dim \Theta _{x,X}=\dim X$ for all $x\in U$ and $\dim  \Theta _{x,X}>\dim X$ for all $x\not \in U$. 
\item Define local parameters, dimension of a local ring, regular local ring and Taylor series expansion.
\item Let $X$ be a quasi-projective variety and $x$ a nonsingular point of $X$. Show that $X$ is irreducibe at $x$.
\item Explain why the set of singular points on a quasi projective variety is closed.
\item Explain why if $x\in X$ is a nonsingular point, then $X$ is locally a complete intersection at $x$.
\item Define local equations of a subvariety $Y\subset X$ at a point $x\in Y$.
\item Explain why a rational map from a non-singular quasi projective variey to projective space is regular in codimension $1$.
\end{enumerate}
\centerline{Homework 10}

\begin{enumerate}
\item 
Define the blow up $\pi:X\to \mathbb P ^n$ centered at the point $p=(1:0:...:0)$ and show that $X$ is smooth and irreducible, $\pi ^{-1}(p)\cong \mathbb P ^{n-1}$ and $X\setminus \pi ^{-1}(p)\cong \mathbb P ^n \setminus p$.

\item Define the normalization of an irreducible quasi projectve variety and compute an example where the induced map is not an isomorphism.

\item Show that an irreducible quasi-projective variety $X$ is normal if and only if $\mathcal O _x$ is normal for all points $x\in X$.

\item Show that smooth varieties are normal.

\item Show that a normal variety is non-singular in codimension $1$.
\end{enumerate}
\end{document}
\end
